
\section{Proposta: un criterio di campionamento consapevole della criticit\`a}

I risultati dell'Esperimento 2 individuano un limite preciso e circoscritto: il criterio di
campionamento proporzionale opera sull'identificativo di traccia, e poich\'e tutti i thread di
un'esecuzione condividono tale identificativo, la decisione ricade sull'intera esecuzione. Non
esiste, con i componenti disponibili, un modo di esprimere una politica del tipo ``conserva sempre
le tracce dei task critici e riduci quelle dei task best-effort''.

In questa sezione delineiamo come tale limite possa essere superato. L'esposizione resta a livello
di progetto: si indicano i punti del sistema su cui intervenire e la logica da introdurvi, senza
scendere nel dettaglio implementativo.

\subsection{Il vincolo da cui partire}

L'analisi del funzionamento della libreria evidenzia due elementi che condizionano qualunque
soluzione.

\textbf{Il momento della decisione.} Il componente che decide se registrare una traccia viene
interpellato \emph{all'atto della creazione} dell'elemento, e riceve il nome dell'elemento e i soli
attributi passati contestualmente alla creazione. Nel codice a disposizione, gli attributi che
descrivono la politica di scheduling e la priorit\`a --- cio\`e proprio le informazioni che
identificano la criticit\`a di un task --- vengono impostati \emph{successivamente}. Un criterio
personalizzato, allo stato attuale, non li vedrebbe.

\textbf{La struttura delle tracce.} Ogni thread viene creato come discendente dell'elemento radice
che rappresenta l'esecuzione. Questa scelta \`e ci\`o che rende l'intera esecuzione una traccia sola,
ed \`e la causa diretta del comportamento osservato in Figura~\ref{fig:01}.

\subsection{Due interventi complementari}

\textbf{Primo intervento: rendere la criticit\`a visibile al momento della decisione.} Sono
possibili due strade, di costo diverso. La meno invasiva consiste nell'adottare una convenzione sui
nomi dei task, che il criterio personalizzato pu\`o esaminare: il nome dell'elemento \`e infatti
gi\`a disponibile al momento della decisione. Le configurazioni della campagna adottano gi\`a nomi
distinti per i task critici e per quelli di disturbo, per cui questa strada non richiederebbe
modifiche al codice esistente. La strada pi\`u pulita, ma che comporta una piccola modifica nel
punto in cui gli elementi dei thread vengono creati, consiste nel passare la criticit\`a come
attributo contestualmente alla creazione, anzich\'e subito dopo. \`E preferibile in prospettiva,
perch\'e non lega la politica a una convenzione sui nomi.

\textbf{Secondo intervento: separare le tracce per criticit\`a.} \`E il punto essenziale, e senza
di esso il primo non basta. Finch\'e tutti i thread appartengono alla stessa traccia, qualunque
criterio proporzionale continuer\`a a valutare un unico identificativo, e le decisioni sui diversi
task resteranno rigidamente accoppiate. Occorre che \emph{ciascun thread avvii una propria traccia},
cos\`i che ogni criticit\`a possa essere campionata in modo indipendente.

La relazione causale con l'esecuzione principale, che ha valore diagnostico e non va perduta, pu\`o
essere mantenuta sostituendo il legame di discendenza con un \emph{collegamento} fra tracce, che la
libreria prevede proprio per correlare tracce distinte. Il risultato \`e che le tracce restano
navigabili come oggi, ma diventano campionabili separatamente.

\subsection{Dove intervenire}

Gli interventi si concentrano in tre punti, tutti circoscritti.

\begin{enumerate}
\item \textbf{Un nuovo componente di campionamento}, da introdurre come file a s\'e stante. Ne
implementerebbe l'interfaccia prevista dalla libreria, esponendo una decisione basata sulla
criticit\`a riconosciuta: conservare sempre gli elementi dei task critici, applicare una
proporzione configurabile agli altri, ed ereditare la decisione del genitore per gli elementi
discendenti, cos\`i che una traccia campionata resti completa.

\item \textbf{La funzione di inizializzazione del tracciamento}, dove il fornitore di tracce viene
costruito e dove viene scelto il criterio di campionamento. \`E il punto in cui il nuovo componente
verrebbe registrato, affiancandolo a quelli esistenti secondo lo stesso schema gi\`a utilizzato per
gli altri parametri di configurazione, in modo che la scelta resti selezionabile e non sostitutiva.

\item \textbf{Il punto di creazione degli elementi associati ai thread}, per i due aspetti discussi
sopra: rendere disponibile la criticit\`a al momento della decisione e avviare una traccia distinta
per thread mantenendo il collegamento con l'esecuzione principale.
\end{enumerate}

\subsection{Come verificarne l'efficacia}

La proposta \`e verificabile con la stessa impostazione sperimentale gi\`a costruita, il che ne
rende la valutazione immediata. Sarebbe sufficiente ripetere l'Esperimento 2 sostituendo il criterio
di campionamento, e osservare la Figura~\ref{fig:01}: con un criterio consapevole della criticit\`a,
i punti dovrebbero disporsi \textbf{nella banda intermedia}, oggi vuota, con gli elementi del task
critico sempre presenti e quelli dei task di disturbo ridotti secondo la proporzione impostata.
La stessa figura che oggi documenta il limite diventerebbe quindi la verifica del suo superamento.

Andrebbe inoltre misurato, con l'impostazione dell'Esperimento 1, il costo del nuovo criterio: la
decisione viene presa a ogni creazione di elemento, e un criterio che esamini attributi o nomi \`e
pi\`u oneroso di uno che valuti un solo identificativo. L'ordine di grandezza atteso \`e tale da non
incidere in modo apprezzabile, ma trattandosi di un sistema real-time la verifica \`e necessaria e
non pu\`o essere data per acquisita.

\subsection{Osservazione conclusiva}

Vale la pena rilevare che i due risultati principali di questo studio indicano due margini di
intervento di natura diversa. Il costo di \emph{costruire} gli elementi di traccia --- circa 14\us{}
per iterazione --- \`e sostenuto necessariamente dal thread che li genera e non \`e eliminabile per
via architetturale. Il costo di \emph{consegnarli}, che come si \`e visto \`e di gran lunga
superiore, \`e invece \emph{collocabile}: la modalit\`a differita lo sposta su un thread di
servizio, sottraendolo al percorso critico. La proposta qui delineata agisce su un terzo asse, il
\emph{volume} della telemetria, rendendone la riduzione selettiva anzich\'e indiscriminata.

Le tre leve sono indipendenti e componibili. Per un sistema a criticit\`a mista, le evidenze
raccolte suggeriscono di adottare la consegna differita, di non affidare al campionamento
proporzionale la protezione dei task critici, e di introdurre un criterio in grado di distinguerli.

\end{document}
